% TOP_PROBLEM: {"id": 30002231, "title": "Nuclear Dimension after Jiang–Su Stabilization", "category_id": 4, "category": {"id": 4, "name": "algebra", "display_name": "Algebra", "description": "Group theory, ring theory, field theory, and algebraic structures.", "slug": "algebra", "order_index": 4, "created_at": "2026-07-31T15:26:25.671Z"}, "status": "open"}
% FIRST_POSTED: 2026-10-06
% ATTEMPT_STATUS: unresolved
% OPENAI_CLAIM: {"schema_version":1,"source_repo":"https://github.com/openai/math","source_commit":"adc7f1241b42e322a6451854ab7e4b4c146bf78a","release_date":"2026-10-06","match":"partial","resolution":"partial","independently_reviewed":false,"second_pass":"n/a","families":[{"family":"291","title":"Cuntz comparison, nuclear dimension, and equivariant Jiang–Su stability","subject":"Operator algebras","manuscripts":[{"title":"Nuclear dimension and Jiang–Su stability without elementary subquotients","pdf_path":"preprints/Nuclear-dimension-and-Jiang-Su-stability-without-elementary-subquotients-September-23-2026/paper.pdf","readme_path":"preprints/Nuclear-dimension-and-Jiang-Su-stability-without-elementary-subquotients-September-23-2026/README.md","date":"2026-09-23","theorem_ref":"Theorem 1.3 (intro:main)"}],"lean":{"doc_path":"lean/docs/291.md","comparators":[],"declarations":[],"covers_main_theorem":false},"reasoning_trace":null}]}
\documentclass[11pt]{article}
\usepackage[margin=25mm]{geometry}
\usepackage{fontspec}
\setmainfont{Latin Modern Roman}
\usepackage{amsmath,amssymb,mathtools}
\usepackage{unicode-math}
\setmathfont{Latin Modern Math}
\usepackage{enumitem}
\usepackage{xurl,hyperref,bookmark}
\hypersetup{colorlinks=true,urlcolor=blue}
\setcounter{secnumdepth}{0}
\setlength{\parindent}{0pt}
\setlength{\parskip}{5pt}
\setlength{\emergencystretch}{3em}
\begin{document}
\section*{Nuclear Dimension after Jiang–Su Stabilization}

\subsection{OpenAI's claim}
Finite nuclear dimension is claimed to characterize Jiang--Su stability for separable nuclear algebras lacking elementary ideal subquotients; moreover, Jiang--Su stabilization always lowers nuclear dimension to at most one. The cited manuscript argues for that result; its relation to this problem is specified below.

\subsection{Outline of the argument}
Manuscript 1 (family 291, 2026-09-23): Localized order-zero approximations control ranks even when traces are infinite on some ideals. A supporting map and comparison yield two-color approximations that survive stabilization and hereditary passage. Conversely, Gaussian sign tensors and pure-state excision construct two full orthogonal central elements from the absence of characters, supplying the absorption criterion.

\subsection{Scope relative to this problem}
Manuscript 1 (family 291, 2026-09-23): The bound dim\_nuc(A\(\otimes\)Z)\(\leq\)1 proves the first branch for separable nuclear A; the record also asks decomposition rank of the C(X) inclusion and permits nonseparable nuclear algebras.

\subsection{Formal verification}
For family 291, the scope page documents the following selections. The selection gives uniform Gamma from real rank zero of the designated tracial ultrapower. The documented comparison scope also treats Jiang--Su absorption and Gamma under strict comparison, with additional trace and target-rank conditions in the projectionless case. Manuscript 1 is not among the papers selected by that page, so the family link does not formalize its main conclusion. It does not establish formal coverage of the complete problem clause.

This site has not executed the Lean files or checked their dependencies. These descriptions report the scope of the linked declarations, rather than a new verification of the proofs.

\subsection{Status caveat}
OpenAI attributes these results to an internal model. Its release notice cautions that results lacking full formalization can contain mistakes. This site has not independently assessed the proofs; a solved label records the scope of the reported claim.

\subsection{Sources}
\begin{itemize}
\item OpenAI, \emph{Nuclear dimension and Jiang–Su stability without elementary subquotients}, 2026-09-23; citation key \texttt{OAI:Nuclear-dimension-and-Jiang-Su-stability-without-elementary-subquotients-September-23-2026}. \href{https://github.com/openai/math/blob/adc7f1241b42e322a6451854ab7e4b4c146bf78a/preprints/Nuclear-dimension-and-Jiang-Su-stability-without-elementary-subquotients-September-23-2026/paper.pdf}{PDF (GitHub)}; \href{https://github.com/openai/math/raw/adc7f1241b42e322a6451854ab7e4b4c146bf78a/preprints/Nuclear-dimension-and-Jiang-Su-stability-without-elementary-subquotients-September-23-2026/paper.pdf}{PDF (direct)}; \href{https://github.com/openai/math/blob/adc7f1241b42e322a6451854ab7e4b4c146bf78a/preprints/Nuclear-dimension-and-Jiang-Su-stability-without-elementary-subquotients-September-23-2026/README.md}{Manuscript page and citation}; \href{https://github.com/openai/math/blob/main/preprints/Nuclear-dimension-and-Jiang-Su-stability-without-elementary-subquotients-September-23-2026/paper.pdf}{Latest version}.
\item Family 291: \href{https://github.com/openai/math/blob/adc7f1241b42e322a6451854ab7e4b4c146bf78a/lean/docs/291.md}{Lean scope}.
\item \href{https://github.com/openai/math/blob/adc7f1241b42e322a6451854ab7e4b4c146bf78a/CONTENTS.md}{OpenAI catalogue entry}.
\end{itemize}
\end{document}
